% ============================================================================
% Chapter 13: Multi-Objective Optimization
% Book source: part1-linear-programming/ch09-multi-objective/
%              multi-objective-optimization_updated.tex
% Exercise numbering synced to \begin{ex} order as of 2026-07-13.
% All computations verified numerically (PuLP 3.3.2 / CBC where LPs appear).
% ============================================================================

\section{Chapter 13: Multi-Objective Optimization}

This chapter has ten exercises. The warm-ups (13.1--13.3) cover the
definition of Pareto optimality, dominance checking in a small table, and
weighted-sum scoring. The core problems (13.4--13.8) work the chapter's main
methods on concrete data: frontier identification, weighted sums, the
$\varepsilon$-constraint method (on a table, on two small LPs, and by hand),
and trade-off analysis. The concept problems (13.9--13.10, with 13.10 as the
challenge) probe the known blind spot of weighted sums on non-convex
frontiers and compute the exact weight intervals that select each vertex in
the chapter's furniture example. Every numerical claim below was verified by
direct computation; the LP-based exercises (13.6 and the furniture sweeps)
were additionally re-solved with PuLP/CBC, and outputs shown are actual. The
book's Selected Solutions for 13.1, 13.4, 13.7, and 13.10 were checked and
found correct; they are reused and expanded here.

% ----------------------------------------------------------------------------
\exsol{13.1}{Defining Pareto Optimality}{1}

For a problem minimizing $f_1$ and $f_2$ over a feasible set $P$, a point
$\x \in P$ is \emph{Pareto optimal} if there is no $\overline{\x} \in P$ with
$f_i(\overline{\x}) \leq f_i(\x)$ for both $i = 1, 2$ and
$f_i(\overline{\x}) < f_i(\x)$ for at least one $i$. In words: no other
feasible point is at least as good in both objectives and strictly better in
one. (Equivalently, adapt the book's maximization definition by flipping the
inequalities.) The Pareto frontier is the set of all such points.

% ----------------------------------------------------------------------------
\exsol{13.2}{Spotting Dominated Designs}{1}

The data (both objectives minimized):
$G = (2,9)$, $H = (4,7)$, $I = (5,8)$, $J = (7,3)$, $K = (6,5)$.

\begin{enumerate}
\item Only contract $I$ is dominated: $H = (4,7)$ beats $I = (5,8)$ in both
cost and delivery time. Checking every other pair, one contract is always
cheaper while the other is faster, so no other domination occurs (e.g.,
$K = (6,5)$ vs.\ $J = (7,3)$: $K$ is cheaper, $J$ is faster).
\item In order of increasing $f_1$, the Pareto frontier is
\[
G = (2,9), \quad H = (4,7), \quad K = (6,5), \quad J = (7,3).
\]
Note the $f_2$ values strictly decrease along this list, as they must along a
frontier of a two-objective minimization problem.
\item True. If $G$ is the \emph{unique} minimizer of $f_1$, every other
contract has strictly larger $f_1$, so no contract can be at least as good as
$G$ in both objectives; hence nothing dominates $G$. (Uniqueness matters: if
two contracts tie for the minimum of $f_1$, the one with larger $f_2$ is
dominated by the other.)
\end{enumerate}

% ----------------------------------------------------------------------------
\exsol{13.3}{Computing Weighted-Sum Scores}{1}

\begin{enumerate}
\item With $(w_1, w_2) = (0.6, 0.4)$, $S = 0.6 f_1 + 0.4 f_2$:
\begin{center}
\begin{tabular}{c|ccccc}
Contract & G & H & I & J & K \\\hline
$S$ & $\mathbf{4.8}$ & $5.2$ & $6.2$ & $5.4$ & $5.6$ \\
\end{tabular}
\end{center}
For example, $S_G = 0.6(2) + 0.4(9) = 1.2 + 3.6 = 4.8$. The winner is
contract $G$.
\item With $(w_1, w_2) = (0.2, 0.8)$:
\begin{center}
\begin{tabular}{c|ccccc}
Contract & G & H & I & J & K \\\hline
$S$ & $7.6$ & $6.4$ & $7.4$ & $\mathbf{3.8}$ & $5.2$ \\
\end{tabular}
\end{center}
The winner is now contract $J$. The winner changes because the weights encode
the decision maker's exchange rate between the objectives: with most of the
weight on delivery time ($w_2 = 0.8$), the contract with by far the best
delivery time ($J$, at 3 days) overcomes its high cost, whereas with the
weight mostly on cost, the cheapest contract ($G$) wins. Geometrically, the
two weight vectors correspond to supporting lines of different slopes, which
touch the frontier at different points.
\item $S_I - S_H = w_1(5 - 4) + w_2(8 - 7) = w_1 + w_2 > 0$ whenever
$w_1, w_2 > 0$. So $I$ scores strictly worse than $H$ for every positive
weight pair, and in general a dominated point loses to any point dominating
it: the score difference is a nonnegative combination of the coordinatewise
gaps, at least one of which is strictly positive.
\end{enumerate}

% ----------------------------------------------------------------------------
\exsol{13.4}{Pareto Frontier, Weighted Sum, and $\varepsilon$-Constraint}{2}

The candidates (both objectives minimized):
$A = (6,5)$, $B = (4,8)$, $C = (7,4)$, $D = (3,6)$, $E = (5,6)$,
$F = (10,3)$.

\begin{enumerate}
\item Compare pairs for domination. $D = (3,6)$ dominates $B = (4,8)$, and
$E = (5,6)$ is dominated by $D$ ($3 \leq 5$, $6 \leq 6$, with the first
strict). $A$, $C$, $F$ are not dominated by anything. The Pareto frontier is
$\{D, A, C, F\}$ with objective vectors $(3,6)$, $(6,5)$, $(7,4)$, $(10,3)$.
\item A design is a weighted-sum optimum for some $\alpha$ exactly when it
lies on the lower-left convex hull of the frontier points. Plotting $(3,6)$,
$(6,5)$, $(7,4)$, $(10,3)$: the hull is formed by $D$, $C$, and $F$, while
$A$ lies above it. Indeed, $D$ (best $f_1$) and $F$ (best $f_2$) are always
obtainable at extreme weights; $C$ lies on the hull between them (the segment
$D \to F$ has slope $-3/7$, and at $f_1 = 7$ it sits at
$6 - 4 \cdot \tfrac{3}{7} \approx 4.29 > 4$, so $C$ is below that segment,
hence on the hull); but $A = (6,5)$ sits above the segment from $D$ to $C$
(that segment at $f_1 = 6$ has value $6 - \tfrac{2}{4}(6-3) = 4.5 < 5$). So
$D$, $C$, $F$ are obtainable by weighted sums and $A$ is not.
\item Choose $\varepsilon = 5$: minimizing $f_1$ subject to $f_2 \leq 5$
leaves candidates $A = (6,5)$, $C = (7,4)$, $F = (10,3)$, and selects
$A = (6,5)$, exactly the frontier point unreachable by any weighted sum.
(Any $\varepsilon \in [5, 6)$ works.)
\item Sketch description: plot the six points in the $(f_1, f_2)$ plane.
Mark $B$ and $E$ as dominated ($B$ is up and to the right of $D$; $E$ is
directly right of $D$). The frontier is the ``staircase'' of points $D, A,
C, F$ running from upper-left to lower-right; improvement arrows point left
(lower cost) and down (lower emissions), toward the empty lower-left region.
A correct sketch shows that no frontier point has another point in its
lower-left quadrant.
\end{enumerate}

% ----------------------------------------------------------------------------
\exsol{13.5}{Weighted Sum Method}{2}

In minimization form: $P1 = (f_1, f_2) = (2, -7)$, $P2 = (5, -9)$,
$P3 = (3, -6)$.

\begin{enumerate}
\item $P1$ and $P2$ are Pareto optimal; $P3$ is not. $P1$ dominates $P3$: it
is cheaper ($2 < 3$) \emph{and} has the higher accessibility score
($7 > 6$). Between $P1$ and $P2$ there is a genuine trade-off ($P1$ cheaper,
$P2$ more accessible), so neither dominates the other.
\item Scores $S = 0.5 f_1 + 0.5 f_2$:
\begin{align*}
S_{P1} &= 0.5(2) + 0.5(-7) = -2.5, \\
S_{P2} &= 0.5(5) + 0.5(-9) = -2.0, \\
S_{P3} &= 0.5(3) + 0.5(-6) = -1.5.
\end{align*}
The lowest score wins: project $P1$ is selected.
\item No. Since $P1$ dominates $P3$,
$S_{P3} - S_{P1} = w_1(3-2) + w_2(-6+7) = w_1 + w_2 > 0$ for all
$w_1, w_2 > 0$, so $P3$ always scores strictly worse than $P1$. Even the
degenerate boundary weights fail to select $P3$: with $w = (1,0)$ (cost
only) $P1$ wins, and with $w = (0,1)$ (accessibility only) $P2$ wins. A
dominated alternative can never be chosen by the weighted-sum method with
positive weights.
\end{enumerate}

% ----------------------------------------------------------------------------
\exsol{13.6}{Epsilon-Constraint Method}{2}

\begin{enumerate}
\item Keeping $f_1$ as the objective and bounding $f_2$:
\begin{align*}
\min \quad & x_1 + 2x_2 \\
\st & 3x_1 + x_2 \leq \varepsilon, \\
& x_1 + x_2 \geq 4, \\
& 0 \leq x_1 \leq 5, \quad 0 \leq x_2 \leq 5.
\end{align*}
\item For $\varepsilon = 10$: the optimum is $(x_1, x_2) = (3, 1)$ with
$f_1 = 5$ and $f_2 = 10$. For $\varepsilon = 6$: the optimum is $(1, 3)$
with $f_1 = 7$ and $f_2 = 6$. In both cases the optimum is the intersection
of $3x_1 + x_2 = \varepsilon$ with $x_1 + x_2 = 4$ (along $x_1 + x_2 = 4$ we
have $f_1 = 8 - x_1$, so the LP pushes $x_1$ as high as the $\varepsilon$
bound allows: $2x_1 = \varepsilon - 4$). Both solutions are Pareto optimal:
each is the unique minimizer of $f_1$ given $f_2 \leq \varepsilon$, and the
$\varepsilon$ bound is active, so no feasible point improves either objective
without worsening the other. Verified with PuLP:

\begin{verbatim}
import pulp

for eps in [10, 6]:
    prob = pulp.LpProblem("biobj", pulp.LpMinimize)
    x1 = pulp.LpVariable("x1", lowBound=0, upBound=5)
    x2 = pulp.LpVariable("x2", lowBound=0, upBound=5)
    prob += x1 + 2*x2                 # f1 (primary objective)
    prob += x1 + x2 >= 4
    prob += 3*x1 + x2 <= eps          # f2 <= epsilon
    prob.solve(pulp.PULP_CBC_CMD(msg=0))
    print(f"eps={eps}: x=({x1.value()}, {x2.value()}), "
          f"f1={x1.value() + 2*x2.value()}, "
          f"f2={3*x1.value() + x2.value()}")
\end{verbatim}

Actual output:
\begin{verbatim}
eps=10: x=(3.0, 1.0), f1=5.0, f2=10.0
eps=6: x=(1.0, 3.0), f1=7.0, f2=6.0
\end{verbatim}

\item Minimizing $f_2$ alone gives $f_2^* = 4$ at $(0, 4)$; minimizing $f_1$
alone gives $f_1 = 4$ at $(4, 0)$, where $f_2 = 12$. So: for
$\varepsilon \geq 12$ the bound is inactive and the solution stays at
$(4,0)$ with $(f_1, f_2) = (4, 12)$. As $\varepsilon$ decreases from $12$
toward $f_2^* = 4$, the optimum slides along the edge $x_1 + x_2 = 4$ from
$(4,0)$ to $(0,4)$, with $x_1 = (\varepsilon - 4)/2$ and
$f_1 = 8 - (\varepsilon - 4)/2$ increasing linearly from $4$ to $8$: each
unit of tightening in $f_2$ costs exactly $\tfrac12$ unit of $f_1$. The
sweep traces the entire Pareto frontier, the segment from $(4,12)$ to
$(8,4)$ in objective space. For $\varepsilon < 4$ the LP is infeasible, so
the sweep stops at $f_2^*$.
\end{enumerate}

% ----------------------------------------------------------------------------
\exsol{13.7}{Epsilon-Constraint by Hand}{2}

\begin{enumerate}
\item At $(0,0)$: $(f_1, f_2) = (0,0)$. At $(6,0)$: $(18, 6)$. At $(0,6)$:
$(6, 18)$. So $(6,0)$ maximizes $f_1$ alone and $(0,6)$ maximizes $f_2$
alone.
\item For $\varepsilon > 6$ the constraint $x_1 + 3x_2 \geq \varepsilon$ is
binding at the optimum, and the solution lies where it meets the edge
$x_1 + x_2 = 6$: solving the two equations gives $x_2 = (\varepsilon - 6)/2$
and $x_1 = (18 - \varepsilon)/2$. (For $\varepsilon = 6$ the vertex $(6,0)$
itself is optimal.) Tabulating:
\begin{center}
\begin{tabular}{c|cccc}
$\varepsilon$ & $x_1$ & $x_2$ & $f_1$ & $f_2$ \\\hline
6  & 6 & 0 & 18 & 6 \\
10 & 4 & 2 & 14 & 10 \\
14 & 2 & 4 & 10 & 14 \\
18 & 0 & 6 & 6  & 18 \\
\end{tabular}
\end{center}
Each row is easily confirmed graphically: the sliced region is the triangle
cut by the line $x_1 + 3x_2 = \varepsilon$, and the objective $3x_1 + x_2$
increases toward the lower-right, so the optimum sits at the corner of the
slice on the hypotenuse.
\item Each tabulated solution has $f_2 = \varepsilon$ exactly and the largest
possible $f_1$ given that bound, so each is Pareto optimal. The Pareto
frontier in decision space is the edge $x_1 + x_2 = 6$ from $(6,0)$ to
$(0,6)$: parametrizing by $x_2 = t$ gives $f_1 = 18 - 2t$ and
$f_2 = 6 + 2t$, so the trade-off rate is one for one ($\Delta f_1 / \Delta
f_2 = -1$): each unit gained in $f_2$ costs exactly one unit of $f_1$.
\item The maximum of $f_2$ over $X$ is $18$ (at $(0,6)$), so for
$\varepsilon > 18$ the constrained LP is \emph{infeasible}; the sweep should
stop at $\varepsilon = 18$. (A practical lesson for implementations: first
compute each objective's individual optimum to set the sweep range.)
\end{enumerate}

% ----------------------------------------------------------------------------
\exsol{13.8}{Trade-Off Analysis}{2}

\begin{enumerate}
\item Plot description: the four points $(40,20)$, $(45,14)$, $(55,10)$,
$(70,8)$ in the (cost, emissions) plane form a decreasing chain; connecting
them from upper-left to lower-right gives the (discrete, convex) Pareto
frontier. No point lies below-left of another, so all four are on the
frontier; improvement is toward the lower-left.
\item Marginal rates of substitution (emissions reduced per \$1000 of extra
cost):
\begin{center}
\begin{tabular}{l|ccc}
Step & $\Delta f_1$ & $\Delta f_2$ & tons per \$1000 \\\hline
S1 $\to$ S2 & $+5$  & $-6$ & $6/5 = 1.20$ \\
S2 $\to$ S3 & $+10$ & $-4$ & $4/10 = 0.40$ \\
S3 $\to$ S4 & $+15$ & $-2$ & $2/15 \approx 0.13$ \\
\end{tabular}
\end{center}
The step S1 $\to$ S2 gives by far the largest emissions reduction per dollar
(1.2 tons per \$1000). The steadily deteriorating ratios ($1.20 \to 0.40 \to
0.13$) mean S2 is the natural elbow (knee point) of this frontier.
\item With a budget of at most $\$50{,}000$, the feasible Pareto solutions
are S1 ($\$40$k) and S2 ($\$45$k); the firm should choose \textbf{S2}: it
fits the budget and has lower emissions (14 vs.\ 20 tons) than S1, and S3
($\$55$k) is unaffordable. (Consistent with part (2), S2 also happens to be
the best value step on the frontier.)
\end{enumerate}

% ----------------------------------------------------------------------------
\exsol{13.9}{What Weighted Sums Cannot See}{2}

The points, both objectives minimized: $P = (0,4)$, $Q = (3,3)$, $R = (4,0)$.

\begin{enumerate}
\item Pairwise: $P$ vs.\ $Q$: $P$ is better in $f_1$ ($0 < 3$), $Q$ better
in $f_2$ ($3 < 4$). $Q$ vs.\ $R$: $Q$ better in $f_1$, $R$ better in $f_2$.
$P$ vs.\ $R$: $P$ better in $f_1$, $R$ better in $f_2$. No domination in any
pair, so all three are Pareto optimal.
\item The weighted scores with weights $(\alpha, 1-\alpha)$:
\[
S_P(\alpha) = 4(1-\alpha) = 4 - 4\alpha, \qquad
S_Q(\alpha) = 3\alpha + 3(1-\alpha) = 3, \qquad
S_R(\alpha) = 4\alpha.
\]
For any $\alpha \in [0,1]$,
$\min\{S_P, S_R\} = \min\{4 - 4\alpha, \, 4\alpha\} \leq 2$, since the two
values sum to $4$ (their minimum is maximized at $\alpha = \tfrac12$, where
both equal $2$). Hence $\min\{S_P, S_R\} \leq 2 < 3 = S_Q$ for every
$\alpha$: either $P$ or $R$ always scores strictly better than $Q$, so $Q$
is never a weighted-sum optimum, not even tied.
\item Geometrically, the weighted-sum method minimizes a linear function
over the points in objective space, so its optima are exactly the points on
the lower-left convex hull, where some supporting line with normal
$(\alpha, 1-\alpha) \geq 0$ touches. Here the hull's lower boundary is the
segment from $P = (0,4)$ to $R = (4,0)$; its midpoint is $(2,2)$, and $Q =
(3,3)$ lies strictly above the segment (the line through $P$ and $R$ is
$f_1 + f_2 = 4$, while $Q$ has $f_1 + f_2 = 6$). Every supporting line
touches only $P$, $R$, or the whole segment, never $Q$. This is precisely
the ``Cons'' entry for the weighted-sum method in Table~13.1: it may miss
Pareto optimal solutions in non-convex regions of the frontier, and this
three-point frontier is non-convex at $Q$.
\item The $\varepsilon$-constraint method recovers $Q$: minimize $f_1$
subject to $f_2 \leq \varepsilon$ with $\varepsilon = 3$. The feasible
points are then $Q$ (with $f_2 = 3$) and $R$ (with $f_2 = 0$), and $Q$ wins
with $f_1 = 3 < 4$. (Goal programming with goals $(3,3)$ and any positive
weights also selects $Q$, whose total deviation is $0$.)
\end{enumerate}

% ----------------------------------------------------------------------------
\exsol{13.10}{Weights That Select Each Vertex}{3}

The book's Selected Solution is correct as printed and is expanded here;
all breakpoints below were re-verified both algebraically and by a PuLP
sweep over $\alpha$.

\begin{enumerate}
\item The weighted objective
$\alpha(8x + 2y) - (1-\alpha)(10x + 2y)$ at each vertex of $P$:
\[
(0,0): \; 0, \qquad
(0,20): \; 80\alpha - 40, \qquad
(10,8): \; 212\alpha - 116, \qquad
(12,0): \; 216\alpha - 120.
\]
(For example, $(10,8)$ has $f_1 = 96$ and $f_2 = 116$, giving
$96\alpha - 116(1 - \alpha) = 212\alpha - 116$.)
\item The optimal value is the upper envelope of these four linear functions
of $\alpha$. Comparing pairs: $(0,20)$ beats $(0,0)$ iff
$80\alpha - 40 \geq 0$, i.e., $\alpha \geq \tfrac12$; $(10,8)$ beats
$(0,20)$ iff $212\alpha - 116 \geq 80\alpha - 40$, i.e., $132\alpha \geq
76$, i.e., $\alpha \geq \tfrac{19}{33} \approx 0.576$; and $(10,8)$ beats
$(12,0)$ iff $4 \geq 4\alpha$, which holds for all $\alpha \leq 1$. Hence
\[
(0,0) \text{ optimal for } \alpha \in [0, \tfrac12], \qquad
(0,20) \text{ for } \alpha \in [\tfrac12, \tfrac{19}{33}], \qquad
(10,8) \text{ for } \alpha \in [\tfrac{19}{33}, 1],
\]
with breakpoints $\alpha = \tfrac12$ and $\alpha = \tfrac{19}{33}$.
A numerical sweep confirms the switches (actual PuLP/CBC output of the
weighted LP): $\alpha = 0.4999$ returns $(0,0)$; $\alpha = 0.5001$ and
$\alpha = 0.5757$ return $(0,20)$; $\alpha = 0.5758$ returns $(10,8)$.
\item At $\alpha = \tfrac12$ the entire edge $x = 0$ is optimal (every plan
$(0, y)$, $0 \leq y \leq 20$, scores $0$); at $\alpha = \tfrac{19}{33}$ the
whole rosewood edge from $(0,20)$ to $(10,8)$ is optimal. At $\alpha = 1$
waste is ignored, so all revenue-optimal plans are optimal: the entire edge
from $(10,8)$ to $(12,0)$, \emph{including the dominated vertex} $(12,0)$,
which earns the same $\$96{,}000$ as $(10,8)$ but wastes $120 > 116$ bdft.
Lesson: a weighted sum with a zero weight can return points that are not
Pareto optimal; any strictly positive weight on waste breaks the tie in
favor of $(10,8)$.
\item With revenue in raw dollars, $f_1 = 8000x + 2000y$, the vertex values
become $0$, $40040\alpha - 40$, $96116\alpha - 116$, and $96120\alpha -
120$, and the same comparisons give breakpoints
\[
\alpha = \tfrac{40}{40040} = \tfrac{1}{1001} \approx 0.001
\qquad\text{and}\qquad
\alpha = \tfrac{76}{56076} = \tfrac{19}{14019} \approx 0.00136 :
\]
virtually every weight in $[0,1]$ now selects $(10,8)$, and the entire
trade-off structure is squeezed into the first $0.14\%$ of the weight range.
When objectives live on wildly different scales, the weights lose their
intended meaning as relative importances, which is why objectives should be
normalized before weighting (Section 13.2).
\end{enumerate}
